\begin{code}[hide]%
\>[0]\AgdaKeyword{module}\AgdaSpace{}%
\AgdaModule{vqupit.sections.composing}\AgdaSpace{}%
\AgdaKeyword{where}\<%
\\
%
\\[\AgdaEmptyExtraSkip]%
\>[0]\AgdaKeyword{open}\AgdaSpace{}%
\AgdaKeyword{import}\AgdaSpace{}%
\AgdaModule{Data.Nat}\AgdaSpace{}%
\AgdaKeyword{using}\AgdaSpace{}%
\AgdaSymbol{(}\AgdaDatatype{ℕ}\AgdaSymbol{)}\<%
\\
\>[0]\AgdaKeyword{open}\AgdaSpace{}%
\AgdaKeyword{import}\AgdaSpace{}%
\AgdaModule{Data.Sum}\AgdaSpace{}%
\AgdaKeyword{using}\AgdaSpace{}%
\AgdaSymbol{(}\AgdaOperator{\AgdaDatatype{\AgdaUnderscore{}⊎\AgdaUnderscore{}}}\AgdaSymbol{)}\<%
\\
\>[0]\AgdaKeyword{import}\AgdaSpace{}%
\AgdaModule{Relation.Binary.PropositionalEquality}\AgdaSpace{}%
\AgdaSymbol{as}\AgdaSpace{}%
\AgdaModule{Eq}\<%
\\
\>[0]\AgdaKeyword{open}\AgdaSpace{}%
\AgdaKeyword{import}\AgdaSpace{}%
\AgdaModule{Notations}\AgdaSpace{}%
\AgdaKeyword{using}\AgdaSpace{}%
\AgdaSymbol{(}\AgdaInductiveConstructor{₁₊}\AgdaSymbol{)}\<%
\\
\>[0]\AgdaKeyword{open}\AgdaSpace{}%
\AgdaKeyword{import}\AgdaSpace{}%
\AgdaModule{Word.Base}\AgdaSpace{}%
\AgdaKeyword{using}\AgdaSpace{}%
\AgdaSymbol{(}\AgdaDatatype{Word}\AgdaSpace{}%
\AgdaSymbol{;}\AgdaSpace{}%
\AgdaFunction{WRel}\AgdaSpace{}%
\AgdaSymbol{;}\AgdaSpace{}%
\AgdaOperator{\AgdaInductiveConstructor{[\AgdaUnderscore{}]ʷ}}\AgdaSpace{}%
\AgdaSymbol{;}\AgdaSpace{}%
\AgdaInductiveConstructor{ε}\AgdaSpace{}%
\AgdaSymbol{;}\AgdaSpace{}%
\AgdaOperator{\AgdaInductiveConstructor{\AgdaUnderscore{}•\AgdaUnderscore{}}}\AgdaSpace{}%
\AgdaSymbol{;}\AgdaSpace{}%
\AgdaOperator{\AgdaFunction{\AgdaUnderscore{}\textasciicircum{}\AgdaUnderscore{}}}\AgdaSpace{}%
\AgdaSymbol{;}\AgdaSpace{}%
\AgdaOperator{\AgdaFunction{\AgdaUnderscore{}ʰ'}}\AgdaSymbol{)}\<%
\\
\>[0]\AgdaKeyword{open}\AgdaSpace{}%
\AgdaKeyword{import}\AgdaSpace{}%
\AgdaModule{Presentation.Definitions}\AgdaSpace{}%
\AgdaKeyword{using}\AgdaSpace{}%
\AgdaSymbol{(}\AgdaOperator{\AgdaRecord{\AgdaUnderscore{}IsPresentationOf\AgdaUnderscore{}}}\AgdaSymbol{)}\<%
\\
\>[0]\AgdaKeyword{open}\AgdaSpace{}%
\AgdaKeyword{import}\AgdaSpace{}%
\AgdaModule{Presentation.Construct.Base}\AgdaSpace{}%
\AgdaKeyword{using}\AgdaSpace{}%
\AgdaSymbol{(}\AgdaOperator{\AgdaDatatype{\AgdaUnderscore{}⋄\AgdaUnderscore{}⋄\AgdaUnderscore{}}}\AgdaSpace{}%
\AgdaSymbol{;}\AgdaSpace{}%
\AgdaDatatype{ConjRelʷ}\AgdaSymbol{)}\<%
\\
%
\\[\AgdaEmptyExtraSkip]%
\>[0]\AgdaKeyword{infixl}\AgdaSpace{}%
\AgdaNumber{9}\AgdaSpace{}%
\AgdaOperator{\AgdaPostulate{\AgdaUnderscore{}↑}}\<%
\\
\>[0]\AgdaKeyword{infix}\AgdaSpace{}%
\AgdaNumber{4}\AgdaSpace{}%
\AgdaOperator{\AgdaDatatype{\AgdaUnderscore{}PRel,\AgdaUnderscore{}===\AgdaUnderscore{}}}\AgdaSpace{}%
\AgdaOperator{\AgdaPostulate{\AgdaUnderscore{}SpRel,\AgdaUnderscore{}===\AgdaUnderscore{}}}\AgdaSpace{}%
\AgdaOperator{\AgdaPostulate{\AgdaUnderscore{}Exact,\AgdaUnderscore{}===\AgdaUnderscore{}}}\<%
\\
%
\\[\AgdaEmptyExtraSkip]%
\>[0]\AgdaKeyword{private}\AgdaSpace{}%
\AgdaKeyword{variable}\<%
\\
\>[0][@{}l@{\AgdaIndent{0}}]%
\>[2]\AgdaGeneralizable{n}\AgdaSpace{}%
\AgdaSymbol{:}\AgdaSpace{}%
\AgdaDatatype{ℕ}\<%
\\
%
\\[\AgdaEmptyExtraSkip]%
\>[0]\AgdaKeyword{module}\AgdaSpace{}%
\AgdaModule{Sym}\AgdaSpace{}%
\AgdaKeyword{where}\<%
\\
\>[0][@{}l@{\AgdaIndent{0}}]%
\>[2]\AgdaKeyword{data}\AgdaSpace{}%
\AgdaDatatype{Gen}\AgdaSpace{}%
\AgdaSymbol{:}\AgdaSpace{}%
\AgdaDatatype{ℕ}\AgdaSpace{}%
\AgdaSymbol{→}\AgdaSpace{}%
\AgdaPrimitive{Set}\AgdaSpace{}%
\AgdaKeyword{where}\<%
\\
\>[2][@{}l@{\AgdaIndent{0}}]%
\>[4]\AgdaInductiveConstructor{H-gen}\AgdaSpace{}%
\AgdaInductiveConstructor{S-gen}\AgdaSpace{}%
\AgdaInductiveConstructor{CZ-gen}\AgdaSpace{}%
\AgdaSymbol{:}\AgdaSpace{}%
\AgdaDatatype{Gen}\AgdaSpace{}%
\AgdaGeneralizable{n}\<%
\\
%
\>[4]\AgdaOperator{\AgdaInductiveConstructor{\AgdaUnderscore{}↥}}\AgdaSpace{}%
\AgdaSymbol{:}\AgdaSpace{}%
\AgdaDatatype{Gen}\AgdaSpace{}%
\AgdaGeneralizable{n}\AgdaSpace{}%
\AgdaSymbol{→}\AgdaSpace{}%
\AgdaDatatype{Gen}\AgdaSpace{}%
\AgdaSymbol{(}\AgdaInductiveConstructor{₁₊}\AgdaSpace{}%
\AgdaGeneralizable{n}\AgdaSymbol{)}\<%
\\
\>[0]\AgdaKeyword{module}\AgdaSpace{}%
\AgdaModule{XZ}\AgdaSpace{}%
\AgdaKeyword{where}\<%
\\
\>[0][@{}l@{\AgdaIndent{0}}]%
\>[2]\AgdaKeyword{data}\AgdaSpace{}%
\AgdaDatatype{Gen}\AgdaSpace{}%
\AgdaSymbol{:}\AgdaSpace{}%
\AgdaDatatype{ℕ}\AgdaSpace{}%
\AgdaSymbol{→}\AgdaSpace{}%
\AgdaPrimitive{Set}\AgdaSpace{}%
\AgdaKeyword{where}\<%
\\
\>[2][@{}l@{\AgdaIndent{0}}]%
\>[4]\AgdaInductiveConstructor{X-gen}\AgdaSpace{}%
\AgdaInductiveConstructor{Z-gen}\AgdaSpace{}%
\AgdaSymbol{:}\AgdaSpace{}%
\AgdaDatatype{Gen}\AgdaSpace{}%
\AgdaGeneralizable{n}\<%
\\
%
\>[4]\AgdaOperator{\AgdaInductiveConstructor{\AgdaUnderscore{}↥}}\AgdaSpace{}%
\AgdaSymbol{:}\AgdaSpace{}%
\AgdaDatatype{Gen}\AgdaSpace{}%
\AgdaGeneralizable{n}\AgdaSpace{}%
\AgdaSymbol{→}\AgdaSpace{}%
\AgdaDatatype{Gen}\AgdaSpace{}%
\AgdaSymbol{(}\AgdaInductiveConstructor{₁₊}\AgdaSpace{}%
\AgdaGeneralizable{n}\AgdaSymbol{)}\<%
\\
%
\>[2]\AgdaKeyword{postulate}\<%
\\
\>[2][@{}l@{\AgdaIndent{0}}]%
\>[4]\AgdaOperator{\AgdaPostulate{\AgdaUnderscore{}QRel,\AgdaUnderscore{}===\AgdaUnderscore{}}}\AgdaSpace{}%
\AgdaSymbol{:}\AgdaSpace{}%
\AgdaDatatype{ℕ}\AgdaSpace{}%
\AgdaSymbol{→}\AgdaSpace{}%
\AgdaPrimitive{Set}\<%
\\
%
\>[4]\AgdaPostulate{presentation}\AgdaSpace{}%
\AgdaSymbol{:}\AgdaSpace{}%
\AgdaPrimitive{Set}\<%
\\
\>[0]\AgdaKeyword{open}\AgdaSpace{}%
\AgdaModule{XZ}\AgdaSpace{}%
\AgdaKeyword{using}\AgdaSpace{}%
\AgdaSymbol{(}\AgdaInductiveConstructor{X-gen}\AgdaSpace{}%
\AgdaSymbol{;}\AgdaSpace{}%
\AgdaInductiveConstructor{Z-gen}\AgdaSpace{}%
\AgdaSymbol{;}\AgdaSpace{}%
\AgdaOperator{\AgdaInductiveConstructor{\AgdaUnderscore{}↥}}\AgdaSymbol{)}\<%
\\
\>[0]\AgdaKeyword{open}\AgdaSpace{}%
\AgdaModule{Sym}\AgdaSpace{}%
\AgdaKeyword{using}\AgdaSpace{}%
\AgdaSymbol{(}\AgdaOperator{\AgdaInductiveConstructor{\AgdaUnderscore{}↥}}\AgdaSymbol{)}\<%
\\
\>[0]\AgdaKeyword{module}\AgdaSpace{}%
\AgdaModule{Sim}\AgdaSpace{}%
\AgdaKeyword{where}\<%
\\
\>[0][@{}l@{\AgdaIndent{0}}]%
\>[2]\AgdaKeyword{postulate}\AgdaSpace{}%
\AgdaOperator{\AgdaPostulate{\AgdaUnderscore{}QRel,\AgdaUnderscore{}===\AgdaUnderscore{}}}\AgdaSpace{}%
\AgdaSymbol{:}\AgdaSpace{}%
\AgdaDatatype{ℕ}\AgdaSpace{}%
\AgdaSymbol{→}\AgdaSpace{}%
\AgdaPrimitive{Set}\<%
\\
\>[0]\AgdaKeyword{module}\AgdaSpace{}%
\AgdaModule{PB}\AgdaSpace{}%
\AgdaKeyword{where}\<%
\\
\>[0][@{}l@{\AgdaIndent{0}}]%
\>[2]\AgdaKeyword{postulate}\AgdaSpace{}%
\AgdaOperator{\AgdaPostulate{\AgdaUnderscore{}≈\AgdaUnderscore{}}}\AgdaSpace{}%
\AgdaSymbol{:}\AgdaSpace{}%
\AgdaPrimitive{Set}\<%
\\
\>[0]\AgdaKeyword{module}\AgdaSpace{}%
\AgdaModule{PapV1}\AgdaSpace{}%
\AgdaKeyword{where}\<%
\\
\>[0][@{}l@{\AgdaIndent{0}}]%
\>[2]\AgdaKeyword{module}\AgdaSpace{}%
\AgdaModule{Clifford-Relations}\AgdaSpace{}%
\AgdaKeyword{where}\<%
\\
\>[2][@{}l@{\AgdaIndent{0}}]%
\>[4]\AgdaKeyword{postulate}\AgdaSpace{}%
\AgdaOperator{\AgdaPostulate{\AgdaUnderscore{}QRel,\AgdaUnderscore{}===\AgdaUnderscore{}}}\AgdaSpace{}%
\AgdaSymbol{:}\AgdaSpace{}%
\AgdaDatatype{ℕ}\AgdaSpace{}%
\AgdaSymbol{→}\AgdaSpace{}%
\AgdaPrimitive{Set}\<%
\\
\>[0]\AgdaKeyword{module}\AgdaSpace{}%
\AgdaModule{V1}\AgdaSpace{}%
\AgdaKeyword{where}\<%
\\
\>[0][@{}l@{\AgdaIndent{0}}]%
\>[2]\AgdaKeyword{module}\AgdaSpace{}%
\AgdaModule{Clifford-Relations}\AgdaSpace{}%
\AgdaKeyword{where}\<%
\\
\>[2][@{}l@{\AgdaIndent{0}}]%
\>[4]\AgdaKeyword{postulate}\AgdaSpace{}%
\AgdaOperator{\AgdaPostulate{\AgdaUnderscore{}QRel,\AgdaUnderscore{}===\AgdaUnderscore{}}}\AgdaSpace{}%
\AgdaSymbol{:}\AgdaSpace{}%
\AgdaDatatype{ℕ}\AgdaSpace{}%
\AgdaSymbol{→}\AgdaSpace{}%
\AgdaPrimitive{Set}\<%
\\
\>[0]\AgdaKeyword{open}\AgdaSpace{}%
\AgdaModule{V1}\<%
\\
\>[0]\AgdaKeyword{module}\AgdaSpace{}%
\AgdaModule{Semidirect}\AgdaSpace{}%
\AgdaKeyword{where}\<%
\\
\>[0][@{}l@{\AgdaIndent{0}}]%
\>[2]\AgdaKeyword{postulate}\AgdaSpace{}%
\AgdaPostulate{Pauli⋊Sp}\AgdaSpace{}%
\AgdaSymbol{:}\AgdaSpace{}%
\AgdaDatatype{ℕ}\AgdaSpace{}%
\AgdaSymbol{→}\AgdaSpace{}%
\AgdaPrimitive{Set}\<%
\\
\>[0]\AgdaKeyword{module}\AgdaSpace{}%
\AgdaModule{PE}\AgdaSpace{}%
\AgdaKeyword{where}\<%
\\
\>[0][@{}l@{\AgdaIndent{0}}]%
\>[2]\AgdaKeyword{postulate}\AgdaSpace{}%
\AgdaPostulate{Presented-group}\AgdaSpace{}%
\AgdaSymbol{:}\AgdaSpace{}%
\AgdaDatatype{ℕ}\AgdaSpace{}%
\AgdaSymbol{→}\AgdaSpace{}%
\AgdaPrimitive{Set}\AgdaSpace{}%
\AgdaSymbol{→}\AgdaSpace{}%
\AgdaPrimitive{Set}\<%
\\
%
\\[\AgdaEmptyExtraSkip]%
\>[0]\AgdaKeyword{postulate}\<%
\\
\>[0][@{}l@{\AgdaIndent{0}}]%
\>[2]\AgdaPostulate{proof-elided}\AgdaSpace{}%
\AgdaSymbol{:}\AgdaSpace{}%
\AgdaSymbol{∀}\AgdaSpace{}%
\AgdaSymbol{\{}\AgdaBound{ℓ}\AgdaSymbol{\}}\AgdaSpace{}%
\AgdaSymbol{\{}\AgdaBound{A}\AgdaSpace{}%
\AgdaSymbol{:}\AgdaSpace{}%
\AgdaPrimitive{Set}\AgdaSpace{}%
\AgdaBound{ℓ}\AgdaSymbol{\}}\AgdaSpace{}%
\AgdaSymbol{→}\AgdaSpace{}%
\AgdaBound{A}\<%
\\
%
\>[2]\AgdaPostulate{Gen}\AgdaSpace{}%
\AgdaSymbol{:}\AgdaSpace{}%
\AgdaDatatype{ℕ}\AgdaSpace{}%
\AgdaSymbol{→}\AgdaSpace{}%
\AgdaPrimitive{Set}\<%
\\
%
\>[2]\AgdaPostulate{X}\AgdaSpace{}%
\AgdaPostulate{Z}\AgdaSpace{}%
\AgdaSymbol{:}\AgdaSpace{}%
\AgdaSymbol{∀}\AgdaSpace{}%
\AgdaSymbol{\{}\AgdaBound{n}\AgdaSymbol{\}}\AgdaSpace{}%
\AgdaSymbol{→}\AgdaSpace{}%
\AgdaDatatype{Word}\AgdaSpace{}%
\AgdaSymbol{(}\AgdaPostulate{Gen}\AgdaSpace{}%
\AgdaBound{n}\AgdaSymbol{)}\<%
\\
%
\>[2]\AgdaPostulate{p}\AgdaSpace{}%
\AgdaPostulate{p-1}\AgdaSpace{}%
\AgdaSymbol{:}\AgdaSpace{}%
\AgdaDatatype{ℕ}\<%
\\
%
\>[2]\AgdaPostulate{+ₚ-group}\AgdaSpace{}%
\AgdaPostulate{Pauli⋊Sp}\AgdaSpace{}%
\AgdaSymbol{:}\AgdaSpace{}%
\AgdaDatatype{ℕ}\AgdaSpace{}%
\AgdaSymbol{→}\AgdaSpace{}%
\AgdaPrimitive{Set}\<%
\\
%
\>[2]\AgdaOperator{\AgdaPostulate{\AgdaUnderscore{}SpRel,\AgdaUnderscore{}===\AgdaUnderscore{}}}\AgdaSpace{}%
\AgdaSymbol{:}\AgdaSpace{}%
\AgdaDatatype{ℕ}\AgdaSpace{}%
\AgdaSymbol{→}\AgdaSpace{}%
\AgdaPrimitive{Set}\<%
\\
%
\>[2]\AgdaOperator{\AgdaPostulate{\AgdaUnderscore{}Exact,\AgdaUnderscore{}===\AgdaUnderscore{}}}\AgdaSpace{}%
\AgdaSymbol{:}\AgdaSpace{}%
\AgdaDatatype{ℕ}\AgdaSpace{}%
\AgdaSymbol{→}\AgdaSpace{}%
\AgdaPrimitive{Set}\<%
\\
%
\>[2]\AgdaOperator{\AgdaPostulate{\AgdaUnderscore{}↑}}\AgdaSpace{}%
\AgdaSymbol{:}\AgdaSpace{}%
\AgdaSymbol{∀}\AgdaSpace{}%
\AgdaSymbol{\{}\AgdaBound{n}\AgdaSymbol{\}}\AgdaSpace{}%
\AgdaSymbol{→}\AgdaSpace{}%
\AgdaDatatype{Word}\AgdaSpace{}%
\AgdaSymbol{(}\AgdaPostulate{Gen}\AgdaSpace{}%
\AgdaBound{n}\AgdaSymbol{)}\AgdaSpace{}%
\AgdaSymbol{→}\AgdaSpace{}%
\AgdaDatatype{Word}\AgdaSpace{}%
\AgdaSymbol{(}\AgdaPostulate{Gen}\AgdaSpace{}%
\AgdaSymbol{(}\AgdaInductiveConstructor{₁₊}\AgdaSpace{}%
\AgdaBound{n}\AgdaSymbol{))}\<%
\\
%
\>[2]\AgdaPostulate{complete}\AgdaSpace{}%
\AgdaPostulate{conjw-sem}\AgdaSpace{}%
\AgdaPostulate{denote-eq}\AgdaSpace{}%
\AgdaPostulate{sem}\AgdaSpace{}%
\AgdaSymbol{:}\AgdaSpace{}%
\AgdaPrimitive{Set}\<%
\\
%
\>[2]\AgdaPostulate{exact-generator-data}\AgdaSpace{}%
\AgdaSymbol{:}\AgdaSpace{}%
\AgdaDatatype{ℕ}\AgdaSpace{}%
\AgdaSymbol{→}\AgdaSpace{}%
\AgdaPrimitive{Set}\<%
\end{code}
\section{Composing: Paulis, the Semidirect Product, and Scalars}\label{sec:composing}

From the symplectic presentation to the Clifford presentation there are
four steps: the Pauli factor; the semidirect product, gluing the Pauli
and symplectic presentations along the conjugation action; plumbing,
which simplifies the glued presentation to the paper's rules over $H$,
$S$, $CZ$ alone; and the central extension by $\langle\omega\rangle$,
which restores the scalars.

\subsection{The Pauli factor}\label{sec:pauli}

The projective Pauli group $\PPauli \cong (\Zp \times \Zp)^n$ has an easy
presentation over generators $X$, $Z$ on each wire:

\begin{code}%
\>[0]\AgdaKeyword{data}\AgdaSpace{}%
\AgdaOperator{\AgdaDatatype{\AgdaUnderscore{}PRel,\AgdaUnderscore{}===\AgdaUnderscore{}}}\AgdaSpace{}%
\AgdaSymbol{:}\AgdaSpace{}%
\AgdaSymbol{(}\AgdaBound{n}\AgdaSpace{}%
\AgdaSymbol{:}\AgdaSpace{}%
\AgdaDatatype{ℕ}\AgdaSymbol{)}\AgdaSpace{}%
\AgdaSymbol{→}\AgdaSpace{}%
\AgdaFunction{WRel}\AgdaSpace{}%
\AgdaSymbol{(}\AgdaPostulate{Gen}\AgdaSpace{}%
\AgdaBound{n}\AgdaSymbol{)}\AgdaSpace{}%
\AgdaKeyword{where}\<%
\\
\>[0][@{}l@{\AgdaIndent{0}}]%
\>[2]\AgdaInductiveConstructor{order-X}%
\>[12]\AgdaSymbol{:}\AgdaSpace{}%
\AgdaSymbol{(}\AgdaInductiveConstructor{₁₊}\AgdaSpace{}%
\AgdaGeneralizable{n}\AgdaSymbol{)}\AgdaSpace{}%
\AgdaOperator{\AgdaDatatype{PRel,}}%
\>[28]\AgdaPostulate{X}\AgdaSpace{}%
\AgdaOperator{\AgdaFunction{\textasciicircum{}}}\AgdaSpace{}%
\AgdaPostulate{p}\AgdaSpace{}%
\AgdaOperator{\AgdaDatatype{===}}\AgdaSpace{}%
\AgdaInductiveConstructor{ε}\<%
\\
%
\>[2]\AgdaInductiveConstructor{order-Z}%
\>[12]\AgdaSymbol{:}\AgdaSpace{}%
\AgdaSymbol{(}\AgdaInductiveConstructor{₁₊}\AgdaSpace{}%
\AgdaGeneralizable{n}\AgdaSymbol{)}\AgdaSpace{}%
\AgdaOperator{\AgdaDatatype{PRel,}}%
\>[28]\AgdaPostulate{Z}\AgdaSpace{}%
\AgdaOperator{\AgdaFunction{\textasciicircum{}}}\AgdaSpace{}%
\AgdaPostulate{p}\AgdaSpace{}%
\AgdaOperator{\AgdaDatatype{===}}\AgdaSpace{}%
\AgdaInductiveConstructor{ε}\<%
\\
%
\>[2]\AgdaInductiveConstructor{comm-Z-X}%
\>[12]\AgdaSymbol{:}\AgdaSpace{}%
\AgdaSymbol{(}\AgdaInductiveConstructor{₁₊}\AgdaSpace{}%
\AgdaGeneralizable{n}\AgdaSymbol{)}\AgdaSpace{}%
\AgdaOperator{\AgdaDatatype{PRel,}}%
\>[28]\AgdaPostulate{Z}\AgdaSpace{}%
\AgdaOperator{\AgdaInductiveConstructor{•}}\AgdaSpace{}%
\AgdaPostulate{X}\AgdaSpace{}%
\AgdaOperator{\AgdaDatatype{===}}\AgdaSpace{}%
\AgdaPostulate{X}\AgdaSpace{}%
\AgdaOperator{\AgdaInductiveConstructor{•}}\AgdaSpace{}%
\AgdaPostulate{Z}\<%
\\
%
\\[\AgdaEmptyExtraSkip]%
\>[0]\AgdaFunction{Theorem-Pauli}\AgdaSpace{}%
\AgdaSymbol{:}\AgdaSpace{}%
\AgdaSymbol{∀}\AgdaSpace{}%
\AgdaSymbol{\{}\AgdaBound{n}\AgdaSymbol{\}}\AgdaSpace{}%
\AgdaSymbol{→}\AgdaSpace{}%
\AgdaSymbol{(}\AgdaOperator{\AgdaPostulate{XZ.\AgdaUnderscore{}QRel,\AgdaUnderscore{}===\AgdaUnderscore{}}}\AgdaSpace{}%
\AgdaBound{n}\AgdaSymbol{)}\AgdaSpace{}%
\AgdaOperator{\AgdaRecord{IsPresentationOf}}\AgdaSpace{}%
\AgdaSymbol{(}\AgdaPostulate{+ₚ-group}\AgdaSpace{}%
\AgdaBound{n}\AgdaSymbol{)}\<%
\\
\>[0]\AgdaFunction{Theorem-Pauli}\AgdaSpace{}%
\AgdaSymbol{=}\AgdaSpace{}%
\AgdaPostulate{XZ.presentation}\<%
\end{code}

The semantic group \aid{+ₚ-group}~$n$ is $\aid{Vec}\;(\Zp \times \Zp)\;n$
under componentwise addition --- the very phase space on which $\Sp$
acts in \S\ref{sec:spsem}, which is what will make the semidirect action
definable.  The normal form is $X^{a}Z^{b}$ on each wire, read off the
exponent vector; soundness, uniqueness and \aid{by-normalization} give the
presentation, and the module also exports the soundness and completeness
maps separately, which the next step consumes.  (The library also
presents the same group compositionally, as an $n$-fold direct product
of $\Zp \times \Zp$ with no normal-form work at all; that presentation is
over a different alphabet, and the circuit-alphabet one is the one the
qupit chain needs.)

\subsection{The semidirect product}\label{sec:semidirect}

This is the wreath-product example of \S\ref{sec:snd} with $\Sp$ in
place of $S_n$: the relations of both factors plus twisted commutation
$h\,n\,h^{-1} = \mathit{conj}\,h\,n$.  Here the action of a symplectic
generator on a Pauli generator is the
paper's Lemmas~2.22--2.23, as words --- $H$ swaps $X$ and $Z$ (up to an
inverse), $S$ sends $X$ to $XZ$, $CZ$ couples the two bottom wires, and
a shifted gate acts on a shifted Pauli only:

\begingroup\renewcommand{\AgdaCodeStyle}{\footnotesize}
\begin{code}%
\>[0]\AgdaFunction{conj}\AgdaSpace{}%
\AgdaSymbol{:}\AgdaSpace{}%
\AgdaDatatype{Sym.Gen}\AgdaSpace{}%
\AgdaGeneralizable{n}\AgdaSpace{}%
\AgdaSymbol{→}\AgdaSpace{}%
\AgdaDatatype{XZ.Gen}\AgdaSpace{}%
\AgdaGeneralizable{n}\AgdaSpace{}%
\AgdaSymbol{→}\AgdaSpace{}%
\AgdaDatatype{Word}\AgdaSpace{}%
\AgdaSymbol{(}\AgdaDatatype{XZ.Gen}\AgdaSpace{}%
\AgdaGeneralizable{n}\AgdaSymbol{)}\<%
\\
\>[0]\AgdaFunction{conj}\AgdaSpace{}%
\AgdaInductiveConstructor{Sym.H-gen}%
\>[17]\AgdaInductiveConstructor{X-gen}%
\>[33]\AgdaSymbol{=}\AgdaSpace{}%
\AgdaPostulate{Z}\<%
\\
\>[0]\AgdaFunction{conj}\AgdaSpace{}%
\AgdaInductiveConstructor{Sym.H-gen}%
\>[17]\AgdaInductiveConstructor{Z-gen}%
\>[33]\AgdaSymbol{=}\AgdaSpace{}%
\AgdaPostulate{X}\AgdaSpace{}%
\AgdaOperator{\AgdaFunction{\textasciicircum{}}}\AgdaSpace{}%
\AgdaPostulate{p-1}\<%
\\
\>[0]\AgdaFunction{conj}\AgdaSpace{}%
\AgdaInductiveConstructor{Sym.S-gen}%
\>[17]\AgdaInductiveConstructor{X-gen}%
\>[33]\AgdaSymbol{=}\AgdaSpace{}%
\AgdaPostulate{X}\AgdaSpace{}%
\AgdaOperator{\AgdaInductiveConstructor{•}}\AgdaSpace{}%
\AgdaPostulate{Z}\<%
\\
\>[0]\AgdaFunction{conj}\AgdaSpace{}%
\AgdaInductiveConstructor{Sym.S-gen}%
\>[17]\AgdaInductiveConstructor{Z-gen}%
\>[33]\AgdaSymbol{=}\AgdaSpace{}%
\AgdaPostulate{Z}\<%
\\
\>[0]\AgdaFunction{conj}\AgdaSpace{}%
\AgdaInductiveConstructor{Sym.CZ-gen}%
\>[17]\AgdaInductiveConstructor{X-gen}%
\>[33]\AgdaSymbol{=}\AgdaSpace{}%
\AgdaPostulate{X}\AgdaSpace{}%
\AgdaOperator{\AgdaInductiveConstructor{•}}\AgdaSpace{}%
\AgdaPostulate{Z}\AgdaSpace{}%
\AgdaOperator{\AgdaPostulate{↑}}\<%
\\
\>[0]\AgdaFunction{conj}\AgdaSpace{}%
\AgdaInductiveConstructor{Sym.CZ-gen}%
\>[17]\AgdaInductiveConstructor{Z-gen}%
\>[33]\AgdaSymbol{=}\AgdaSpace{}%
\AgdaPostulate{Z}\<%
\\
\>[0]\AgdaFunction{conj}\AgdaSpace{}%
\AgdaInductiveConstructor{Sym.CZ-gen}%
\>[17]\AgdaSymbol{(}\AgdaInductiveConstructor{X-gen}\AgdaSpace{}%
\AgdaOperator{\AgdaInductiveConstructor{↥}}\AgdaSymbol{)}%
\>[33]\AgdaSymbol{=}\AgdaSpace{}%
\AgdaPostulate{X}\AgdaSpace{}%
\AgdaOperator{\AgdaPostulate{↑}}\AgdaSpace{}%
\AgdaOperator{\AgdaInductiveConstructor{•}}\AgdaSpace{}%
\AgdaPostulate{Z}\<%
\\
\>[0]\AgdaFunction{conj}\AgdaSpace{}%
\AgdaInductiveConstructor{Sym.CZ-gen}%
\>[17]\AgdaSymbol{(}\AgdaInductiveConstructor{Z-gen}\AgdaSpace{}%
\AgdaOperator{\AgdaInductiveConstructor{↥}}\AgdaSymbol{)}%
\>[33]\AgdaSymbol{=}\AgdaSpace{}%
\AgdaPostulate{Z}\AgdaSpace{}%
\AgdaOperator{\AgdaPostulate{↑}}\<%
\\
\>[0]\AgdaFunction{conj}\AgdaSpace{}%
\AgdaInductiveConstructor{Sym.H-gen}%
\>[17]\AgdaSymbol{(}\AgdaBound{xz}\AgdaSpace{}%
\AgdaOperator{\AgdaInductiveConstructor{↥}}\AgdaSymbol{)}%
\>[33]\AgdaSymbol{=}\AgdaSpace{}%
\AgdaOperator{\AgdaInductiveConstructor{[}}\AgdaSpace{}%
\AgdaBound{xz}\AgdaSpace{}%
\AgdaOperator{\AgdaInductiveConstructor{]ʷ}}\AgdaSpace{}%
\AgdaOperator{\AgdaPostulate{↑}}\<%
\\
\>[0]\AgdaFunction{conj}\AgdaSpace{}%
\AgdaInductiveConstructor{Sym.S-gen}%
\>[17]\AgdaSymbol{(}\AgdaBound{xz}\AgdaSpace{}%
\AgdaOperator{\AgdaInductiveConstructor{↥}}\AgdaSymbol{)}%
\>[33]\AgdaSymbol{=}\AgdaSpace{}%
\AgdaOperator{\AgdaInductiveConstructor{[}}\AgdaSpace{}%
\AgdaBound{xz}\AgdaSpace{}%
\AgdaOperator{\AgdaInductiveConstructor{]ʷ}}\AgdaSpace{}%
\AgdaOperator{\AgdaPostulate{↑}}\<%
\\
\>[0]\AgdaFunction{conj}\AgdaSpace{}%
\AgdaInductiveConstructor{Sym.CZ-gen}%
\>[17]\AgdaSymbol{(}\AgdaBound{xz}\AgdaSpace{}%
\AgdaOperator{\AgdaInductiveConstructor{↥}}\AgdaSpace{}%
\AgdaOperator{\AgdaInductiveConstructor{↥}}\AgdaSymbol{)}%
\>[33]\AgdaSymbol{=}\AgdaSpace{}%
\AgdaOperator{\AgdaInductiveConstructor{[}}\AgdaSpace{}%
\AgdaBound{xz}\AgdaSpace{}%
\AgdaOperator{\AgdaInductiveConstructor{]ʷ}}\AgdaSpace{}%
\AgdaOperator{\AgdaPostulate{↑}}\AgdaSpace{}%
\AgdaOperator{\AgdaPostulate{↑}}\<%
\\
\>[0]\AgdaFunction{conj}\AgdaSpace{}%
\AgdaSymbol{(}\AgdaBound{sym}\AgdaSpace{}%
\AgdaOperator{\AgdaInductiveConstructor{↥}}\AgdaSymbol{)}%
\>[17]\AgdaInductiveConstructor{X-gen}%
\>[33]\AgdaSymbol{=}\AgdaSpace{}%
\AgdaPostulate{X}\<%
\\
\>[0]\AgdaFunction{conj}\AgdaSpace{}%
\AgdaSymbol{(}\AgdaBound{sym}\AgdaSpace{}%
\AgdaOperator{\AgdaInductiveConstructor{↥}}\AgdaSymbol{)}%
\>[17]\AgdaInductiveConstructor{Z-gen}%
\>[33]\AgdaSymbol{=}\AgdaSpace{}%
\AgdaPostulate{Z}\<%
\\
\>[0]\AgdaFunction{conj}\AgdaSpace{}%
\AgdaSymbol{(}\AgdaBound{sym}\AgdaSpace{}%
\AgdaOperator{\AgdaInductiveConstructor{↥}}\AgdaSymbol{)}%
\>[17]\AgdaSymbol{(}\AgdaBound{xz}\AgdaSpace{}%
\AgdaOperator{\AgdaInductiveConstructor{↥}}\AgdaSymbol{)}%
\>[33]\AgdaSymbol{=}\AgdaSpace{}%
\AgdaFunction{conj}\AgdaSpace{}%
\AgdaBound{sym}\AgdaSpace{}%
\AgdaBound{xz}\AgdaSpace{}%
\AgdaOperator{\AgdaPostulate{↑}}\<%
\\
%
\\[\AgdaEmptyExtraSkip]%
\>[0]\AgdaFunction{Pauli⋊Symplectic}\AgdaSpace{}%
\AgdaSymbol{:}\AgdaSpace{}%
\AgdaSymbol{(}\AgdaBound{n}\AgdaSpace{}%
\AgdaSymbol{:}\AgdaSpace{}%
\AgdaDatatype{ℕ}\AgdaSymbol{)}\AgdaSpace{}%
\AgdaSymbol{→}\AgdaSpace{}%
\AgdaFunction{WRel}\AgdaSpace{}%
\AgdaSymbol{(}\AgdaDatatype{XZ.Gen}\AgdaSpace{}%
\AgdaBound{n}\AgdaSpace{}%
\AgdaOperator{\AgdaDatatype{⊎}}\AgdaSpace{}%
\AgdaDatatype{Sym.Gen}\AgdaSpace{}%
\AgdaBound{n}\AgdaSymbol{)}\<%
\\
\>[0]\AgdaFunction{Pauli⋊Symplectic}\AgdaSpace{}%
\AgdaBound{n}\AgdaSpace{}%
\AgdaSymbol{=}\AgdaSpace{}%
\AgdaSymbol{(}\AgdaBound{n}\AgdaSpace{}%
\AgdaOperator{\AgdaDatatype{PRel,\AgdaUnderscore{}===\AgdaUnderscore{}}}\AgdaSymbol{)}\AgdaSpace{}%
\AgdaOperator{\AgdaDatatype{⋄}}\AgdaSpace{}%
\AgdaSymbol{(}\AgdaBound{n}\AgdaSpace{}%
\AgdaOperator{\AgdaPostulate{SpRel,\AgdaUnderscore{}===\AgdaUnderscore{}}}\AgdaSymbol{)}\AgdaSpace{}%
\AgdaOperator{\AgdaDatatype{⋄}}\AgdaSpace{}%
\AgdaDatatype{ConjRelʷ}\AgdaSpace{}%
\AgdaFunction{conj}\<%
\end{code}
\endgroup

(module qualifiers on the Pauli side elided; the relation is
\aid{{\_}QRel,{\_}==={\_}} in the source).

The generic theorem of \S\ref{sec:constructions} asks for two things
beyond the factor presentations: that \aid{conj} respects the symplectic
rules in its first argument and the Pauli rules in its second.  Neither
is proved by hand, in contrast to the case splits of \S\ref{sec:snd}.
Both go through the semantics: the action of a symplectic word $c$ on a
Pauli word $w$, computed syntactically by $(\mathit{conj}\,\aid{ʰ'})$,
denotes $\aid{ap}\;\sem{c}\;(\mathit{sem}\;w)$ in the phase space, so if
$c \approx d$ by a symplectic axiom then the two conjugates have equal
denotations by symplectic soundness, and are therefore congruent Pauli
words by Pauli \emph{completeness}:

\begingroup\renewcommand{\AgdaCodeStyle}{\footnotesize}
\begin{code}%
\>[0]\AgdaFunction{respects-Δ}\AgdaSpace{}%
\AgdaSymbol{:}%
\>[350I]\AgdaSymbol{∀}\AgdaSpace{}%
\AgdaSymbol{\{}\AgdaBound{n}\AgdaSymbol{\}}\AgdaSpace{}%
\AgdaSymbol{\{}\AgdaBound{c}\AgdaSpace{}%
\AgdaBound{d}\AgdaSpace{}%
\AgdaSymbol{:}\AgdaSpace{}%
\AgdaDatatype{Word}\AgdaSpace{}%
\AgdaSymbol{(}\AgdaPostulate{Gen}\AgdaSpace{}%
\AgdaBound{n}\AgdaSymbol{)\}}\AgdaSpace{}%
\AgdaSymbol{(}\AgdaBound{u}\AgdaSpace{}%
\AgdaSymbol{:}\AgdaSpace{}%
\AgdaDatatype{Word}\AgdaSpace{}%
\AgdaSymbol{(}\AgdaDatatype{XZ.Gen}\AgdaSpace{}%
\AgdaBound{n}\AgdaSymbol{))}\AgdaSpace{}%
\AgdaSymbol{→}\AgdaSpace{}%
\AgdaOperator{\AgdaPostulate{Sim.\AgdaUnderscore{}QRel,\AgdaUnderscore{}===\AgdaUnderscore{}}}\AgdaSpace{}%
\AgdaBound{n}\AgdaSpace{}%
\AgdaBound{c}\AgdaSpace{}%
\AgdaBound{d}\AgdaSpace{}%
\AgdaSymbol{→}\<%
\\
\>[.][@{}l@{}]\<[350I]%
\>[13]\AgdaOperator{\AgdaPostulate{PB.\AgdaUnderscore{}≈\AgdaUnderscore{}}}\AgdaSpace{}%
\AgdaSymbol{(}\AgdaOperator{\AgdaPostulate{XZ.\AgdaUnderscore{}QRel,\AgdaUnderscore{}===\AgdaUnderscore{}}}\AgdaSpace{}%
\AgdaBound{n}\AgdaSymbol{)}\AgdaSpace{}%
\AgdaSymbol{((}\AgdaFunction{conj}\AgdaSpace{}%
\AgdaOperator{\AgdaFunction{ʰ'}}\AgdaSymbol{)}\AgdaSpace{}%
\AgdaBound{c}\AgdaSpace{}%
\AgdaBound{u}\AgdaSymbol{)}\AgdaSpace{}%
\AgdaSymbol{((}\AgdaFunction{conj}\AgdaSpace{}%
\AgdaOperator{\AgdaFunction{ʰ'}}\AgdaSymbol{)}\AgdaSpace{}%
\AgdaBound{d}\AgdaSpace{}%
\AgdaBound{u}\AgdaSymbol{)}\<%
\\
\>[0]\AgdaFunction{respects-Δ}\AgdaSpace{}%
\AgdaSymbol{\{}\AgdaBound{n}\AgdaSymbol{\}}\AgdaSpace{}%
\AgdaSymbol{\{}\AgdaBound{c}\AgdaSymbol{\}}\AgdaSpace{}%
\AgdaSymbol{\{}\AgdaBound{d}\AgdaSymbol{\}}\AgdaSpace{}%
\AgdaBound{u}\AgdaSpace{}%
\AgdaBound{ax}\AgdaSpace{}%
\AgdaSymbol{=}\<%
\\
\>[0][@{}l@{\AgdaIndent{0}}]%
\>[2]\AgdaPostulate{complete}\AgdaSpace{}%
\AgdaBound{n}\AgdaSpace{}%
\AgdaSymbol{(}\AgdaFunction{Eq.trans}\AgdaSpace{}%
\AgdaSymbol{(}\AgdaPostulate{conjw-sem}\AgdaSpace{}%
\AgdaBound{c}\AgdaSpace{}%
\AgdaBound{u}\AgdaSymbol{)}\AgdaSpace{}%
\AgdaSymbol{(}\AgdaFunction{Eq.trans}\AgdaSpace{}%
\AgdaSymbol{(}\AgdaPostulate{denote-eq}\AgdaSpace{}%
\AgdaBound{ax}\AgdaSpace{}%
\AgdaSymbol{(}\AgdaPostulate{sem}\AgdaSpace{}%
\AgdaBound{u}\AgdaSymbol{))}\AgdaSpace{}%
\AgdaSymbol{(}\AgdaFunction{Eq.sym}\AgdaSpace{}%
\AgdaSymbol{(}\AgdaPostulate{conjw-sem}\AgdaSpace{}%
\AgdaBound{d}\AgdaSpace{}%
\AgdaBound{u}\AgdaSymbol{))))}\<%
\end{code}
\endgroup

\noindent and symmetrically for \aid{respects-Γ}.  The long symplectic
rules --- \aid{M-power}, \aid{semi-M↑CZ}, \aid{selinger-c10} to
\aid{c15} --- are never conjugated by hand.  This is the first place the
factored route pays off in a way the paper's route cannot: the paper's
``Pauli complete'' rules (its Definition~4.7) are precisely these
conjugation rules, and there they must be derived syntactically from
Figure~1 (its Proposition~4.8).  With every premise a theorem, the
presentation \aid{(SemiDirect.{\_}QRel,{\_}==={\_} n) IsPresentationOf
Pauli⋊Sp} is unconditional, where \aid{Pauli⋊Sp} is the theorem's
\aid{G1⋊G2}.

\paragraph{The semantic side.}  \aid{Pauli⋊Sp} is the library's
semidirect product of the two semantic groups, built from an
\aid{Action} record (\S\ref{sec:constructions}) whose action is the
syntactic \aid{conj} transported through the two presentation
isomorphisms.  The \emph{natural} semidirect product, with $\Sp$ acting
on $\Zp^{2n}$ by \aid{ap}, is a different term, and the two are
identified by a lemma $\aid{act≡ap}$ proving that the transported
action is \aid{ap}.  This lemma lives in an \aid{abstract} block: with
it transparent, the file's header records, a two-line lemma about
$Z^{j}$ ``went from 27\,s to out-of-memory at 12\,GB'', because Agda
tried to unfold the transported action through both presentation
isomorphisms.

\subsection{Plumbing: presentation simplification}\label{sec:plumbing}

Two presentations are \emph{equivalent} when the presented groups are
isomorphic --- when each side's axioms are derivable from the other's.
The semidirect presentation has five generators per wire, $X, Z, H, S,
CZ$, and the relations of Table~\ref{tab:plumb}; the paper has three,
$H, S, CZ$, and sixteen relations.  Getting from one to the other is
\emph{generator reduction} --- $X$ and $Z$ become the derived words of
\S\ref{sec:background} --- and \emph{relation reduction} --- the Pauli
rules and the fourteen conjugation rules become derivable.  For
example, $CZ \bullet X{\uparrow} \approx X{\uparrow} \bullet
Z{\downarrow} \bullet CZ$ (\aid{conj CZ-gen (X-gen ↥)}) becomes a
theorem of the Clifford rules.

\begin{table}[h]
\caption{The three presentations of the projective Clifford group and
the two reductions between them.}
\label{tab:plumb}
\centering\scriptsize
\setlength{\tabcolsep}{4pt}
\begin{tabular}{llll}
\toprule
 & \texttt{SemiDirect} & \texttt{Simplified-V1} & \texttt{Paper-V0} / \texttt{Paper-V1} \\
\midrule
generators & $X, Z, H, S, CZ$ & $H, S, CZ$ & $H, S, CZ$ \\
relations  & Pauli (3) + symplectic (15) + conjugation (14) & 19 & 16 / 15 \\
\midrule
\multicolumn{4}{l}{\emph{generator reduction}: $X, Z$ become the derived words; $S \mapsto R = S\cdot Z^{1/2}$}\\
\multicolumn{4}{l}{\emph{relation reduction}: Pauli $+$ conjugation rules become derivable; c10--c15 and the swap calculus interderivable}\\
\bottomrule
\end{tabular}
\end{table}

\paragraph{From the semidirect relation to Simplified-V1.}  The first
reduction is an isomorphism of presentations on \emph{different}
alphabets, given by generator maps in both directions: $f$ sends $X$
and $Z$ to their derived words, $H$ and $CZ$ to themselves, and $S$ to
$R$; $h$ sends $H$ and $CZ$ to themselves and $S$ to $Z^{-1/2} \bullet
S$; both commute with $\uparrow$.  The symplectic $S$ is carried to
$R = S \cdot Z^{1/2}$ and $h$ compensates, because the target rule set,
\texttt{Simplified-V1}, spells the multiplier over $R$, as
\S\ref{sec:background} does: $M_x = R^{x}HR^{x^{-1}}HR^{x}H$ exactly, with no Pauli
correction.  \texttt{Simplified-V1} has nineteen group-specific rules:
the simplified symplectic rules of \S\ref{sec:reduction} with $R$ in
place of $S$ in the multiplier rules and in c10, c11, plus
$(SH)^3 = \varepsilon$, $H^2SH^2S = SH^2SH^2$, and the two rules for
pushing $X$ through $CZ$.  Notably $XZ = ZX$ is \emph{not} among them;
it is derived from $(SH)^3 = \varepsilon$ and \aid{M-power}, by the
paper's own argument (its Proposition~4.8).  The isomorphism theorem
needs the four generator-level facts of the library's
\aid{StarGroupIsomorphism}: $f$ and $h$ respect the axioms
(\aid{f-well-defined}, \aid{h-well-defined}, one clause per axiom, the
former containing the derivations of all three Pauli rules and all
fourteen conjugation rules), and each is a left inverse of the other on
generators.  Composing with the semidirect theorem:

\begin{code}%
\>[0]\AgdaFunction{presentation}\AgdaSpace{}%
\AgdaSymbol{:}\AgdaSpace{}%
\AgdaSymbol{∀}\AgdaSpace{}%
\AgdaSymbol{\{}\AgdaBound{n}\AgdaSymbol{\}}\AgdaSpace{}%
\AgdaSymbol{→}\<%
\\
\>[0][@{}l@{\AgdaIndent{0}}]%
\>[2]\AgdaSymbol{(}\AgdaOperator{\AgdaPostulate{Clifford-Relations.\AgdaUnderscore{}QRel,\AgdaUnderscore{}===\AgdaUnderscore{}}}\AgdaSpace{}%
\AgdaBound{n}\AgdaSymbol{)}\AgdaSpace{}%
\AgdaOperator{\AgdaRecord{IsPresentationOf}}\AgdaSpace{}%
\AgdaSymbol{(}\AgdaPostulate{Semidirect.Pauli⋊Sp}\AgdaSpace{}%
\AgdaBound{n}\AgdaSymbol{)}\<%
\end{code}
\begin{code}[hide]%
\>[0]\AgdaFunction{presentation}\AgdaSpace{}%
\AgdaSymbol{=}\AgdaSpace{}%
\AgdaPostulate{proof-elided}\<%
\end{code}

This is where the mathematics of the composition step lives (about
6\,000 lines): a multiplier calculus, a Pauli-conjugation calculus, the
$CZ$ layer, and a small rewriting tactic that discharges the
bookkeeping the paper's hand derivations spend most of their length
on.

\paragraph{From Simplified-V1 to our sixteen rules.}
\texttt{Paper-V0} is the rule set of \S\ref{sec:background}: the same
alphabet, so the identity on words is the candidate isomorphism.  Ten of
the axioms are shared verbatim; the rest is real work in both
directions.  \texttt{Paper-V0} axiomatises the swap directly and
two-qupit interactions through \aid{blake-c12} and
\aid{semi-CX↑-CZ↓}, and from these its \texttt{Lemmas} (4\,750 lines)
derives Selinger's c10--c15 natively together with the lower-wire
rules, the two Pauli-versus-$CZ$ rules and the swap--$CZ$ commutation;
conversely the seven swap rules are proved in \texttt{Simplified-V1} by
transport along the chain symplectic $\to$ simplified $\to$ semidirect
$\to$ V1 rather than by fresh derivations.  The module exports \aid{v1⇒pap}, which
transports any Simplified-V1 theorem into \texttt{Paper-V0}'s theory;
it is what lets later layers reuse the calculi rather than re-prove
them.

\paragraph{From our sixteen rules to the paper's fifteen: connecting
the two routes.}  The paper's Figure~1 spells the multiplier over
$S$, with an explicit Pauli prefix (its Lemma~2.14): $M_x =
\aid{SHS'}\;x^{-1}\;x$ where $\aid{SHS'}\;a\;b = Z^{(b-1)/2}
X^{(1-a)/2} S^{b} H S^{a} H S^{b} H$, and correspondingly states
\aid{semi-MR} as
$M_g \bullet S^{g^2} \bullet Z^{(g^2-g)/2} = S \bullet M_g$ (its rule
C4), the $Z$-exponent being what is left after the $Z^{1/2}$ hidden in
each $R$ is split off and carried across the multiplier, where it picks
up a factor of $g$.  \texttt{Paper-V1} is Figure~1 modulo scalars in
this spelling: our sixteen rules with \aid{order-H},
\aid{M-power}, \aid{semi-MR} and \aid{semi-M↑CZ} restated over
\aid{SHS'}, and \emph{fifteen} rules rather than sixteen, because
$(SH)^3 = \varepsilon$ is now a theorem.
Under the $S$-spelling, $M_1$ \emph{is} the word $(SH)^3$ --- at $x = 1$
the Pauli prefix vanishes --- and \aid{M-power} at $k = 0$ already
gives $M_1 \approx \varepsilon$; under our $R$-spelling $M_1$ is
$(RH)^3$, a different word, which is why the axiom had to stay there.
The isomorphism with \texttt{Paper-V0} is again the identity on words;
eleven axioms cross directly, and the four multiplier rules need the
bridge between the two spellings, $R^{a}HR^{b}HR^{a}H \approx
\aid{SHS'}\;a\;b$, proved once over an abstract relation in
\texttt{Shared.PauliBase} and transported by \aid{v1⇒pap}.  The result
is the theorem the index module states:

\begin{code}%
\>[0]\AgdaFunction{clifford-presentation}\AgdaSpace{}%
\AgdaSymbol{:}\<%
\\
\>[0][@{}l@{\AgdaIndent{0}}]%
\>[2]\AgdaSymbol{∀}\AgdaSpace{}%
\AgdaBound{n}\AgdaSpace{}%
\AgdaSymbol{→}\AgdaSpace{}%
\AgdaSymbol{(}\AgdaOperator{\AgdaPostulate{PapV1.Clifford-Relations.\AgdaUnderscore{}QRel,\AgdaUnderscore{}===\AgdaUnderscore{}}}\AgdaSpace{}%
\AgdaBound{n}\AgdaSymbol{)}\AgdaSpace{}%
\AgdaOperator{\AgdaRecord{IsPresentationOf}}\AgdaSpace{}%
\AgdaSymbol{(}\AgdaPostulate{Pauli⋊Sp}\AgdaSpace{}%
\AgdaBound{n}\AgdaSymbol{)}\<%
\end{code}
\begin{code}[hide]%
\>[0]\AgdaFunction{clifford-presentation}\AgdaSpace{}%
\AgdaSymbol{=}\AgdaSpace{}%
\AgdaPostulate{proof-elided}\<%
\end{code}

So the paper's Figure~1, read modulo scalars, presents the same group as
the semidirect presentation of \S\ref{sec:semidirect}: the two routes ---
the paper's single rule set with its projective use of relations, and
our factored construction --- provably arrive at the same place.

\subsection{Getting the scalars back}\label{sec:scalars}

The Clifford group is a central extension of the projective Clifford
group by the scalars $\langle\omega\rangle$.  Following the extension
recipe of \S\ref{sec:constructions}, the projective presentation is
modified in four moves: add a generator $\omega$; add the relation
$\omega^{p} = \varepsilon$; add relations that $\omega$ commutes with
every gate; and \emph{correct} each projective relation by the scalar it
actually picks up.  In Agda the recipe is instantiated rather than
spelled out: \aid{extension-presentation} is applied to the cyclic
relation $\langle \omega \mid \omega^{p} = \varepsilon\rangle$, the
projective rules, the constant conjugation $\mathit{conj}\;\_\;\_ =
\omega$, and a correction function that is $\varepsilon$ on every rule
except \aid{order-SH}, where it is $\omega^{(p^2-1)/8}$.  Reading the
generators as the paper's matrices, with $H$ normalised by
$\delta_p$ so that all determinants are one, exactly \emph{one} of the
sixteen rules fails to hold on the nose: $(SH)^3 = \varepsilon$ picks
up $\omega^{(p^2-1)/8}$ (cf.\ the phase in the paper's (T1)), which is
exact division since $8 \mid p^2 - 1$ for odd $p$, and whose residue
modulo $p$ is $-1/8$ --- so the corrected rule reads
\begin{center}
  \fitfig[0.5\textwidth]{sc-order-SH}
\end{center}
where $1/8$ is the inverse of $8$ in $\Zp$.  Every other rule holds
exactly: the three multiplier rules and \aid{order-H} because $R = S
\cdot Z^{1/2}$ has no Pauli part, \aid{blake-c12} because every Pauli in
it is pure-$Z$, and pure-$Z$ Paulis have vanishing symplectic form.  Note that the quotient here is
our sixteen-rule set; correcting the paper's fifteen would give
Figure~1 itself.

\paragraph{The semantic side.}  A central extension of a group $Q$ by
an abelian group $K$ is determined by a 2-cocycle $c : Q \times Q \to
K$, with multiplication $(k, g)\cdot(k', h) = (k + k' + c(g,h),\; gh)$.
The formalisation builds two.  The \emph{structural} one is the Weyl
cocycle $c((P,S),(Q,\_)) = \tfrac12\,\mathrm{sform}(P, \aid{ap}\;S\;Q)$
on the natural semidirect product, whose cocycle law is symplectic
invariance of \aid{sform}, and which
satisfies the Heisenberg law $P \cdot Q = \omega^{\mathrm{sform}(P,Q)}
\, Q \cdot P$ on Paulis --- what distinguishes it from the direct product
$\langle\omega\rangle \times (\PPauli \rtimes \Sp)$.  The
\emph{presented} one, \aid{Presented-group}, has as quotient the word
group of the sixteen rules and as kernel the word group of
$\langle\omega \mid \omega^{p}\rangle$, and its cocycle is generated
from \emph{generator data}: for each gate $a$ and word $v$, the scalar
by which the lifts of $a$ and $v$ fail to compose, together with the
proof that this descends to the quotient and assigns each relator the
same correction on both sides.  The generator data is obtained by pulling
the Weyl cocycle back along the projective presentation theorem, so the
two cocycles agree by construction.  The theorem, again on no hypothesis,
is

\begin{code}%
\>[0]\AgdaFunction{presentation-exact}\AgdaSpace{}%
\AgdaSymbol{:}\AgdaSpace{}%
\AgdaSymbol{∀}\AgdaSpace{}%
\AgdaSymbol{(}\AgdaBound{n}\AgdaSpace{}%
\AgdaSymbol{:}\AgdaSpace{}%
\AgdaDatatype{ℕ}\AgdaSymbol{)}\AgdaSpace{}%
\AgdaSymbol{→}\<%
\\
\>[0][@{}l@{\AgdaIndent{0}}]%
\>[2]\AgdaKeyword{let}\AgdaSpace{}%
\AgdaBound{Clifford-group}\AgdaSpace{}%
\AgdaSymbol{=}\AgdaSpace{}%
\AgdaSymbol{λ}\AgdaSpace{}%
\AgdaBound{n}\AgdaSpace{}%
\AgdaSymbol{→}\AgdaSpace{}%
\AgdaPostulate{PE.Presented-group}\AgdaSpace{}%
\AgdaBound{n}\AgdaSpace{}%
\AgdaSymbol{(}\AgdaPostulate{exact-generator-data}\AgdaSpace{}%
\AgdaBound{n}\AgdaSymbol{)}\<%
\\
%
\>[2]\AgdaKeyword{in}\AgdaSpace{}%
\AgdaSymbol{(}\AgdaBound{n}\AgdaSpace{}%
\AgdaOperator{\AgdaPostulate{Exact,\AgdaUnderscore{}===\AgdaUnderscore{}}}\AgdaSymbol{)}\AgdaSpace{}%
\AgdaOperator{\AgdaRecord{IsPresentationOf}}\AgdaSpace{}%
\AgdaSymbol{(}\AgdaBound{Clifford-group}\AgdaSpace{}%
\AgdaBound{n}\AgdaSymbol{)}\<%
\end{code}
\begin{code}[hide]%
\>[0]\AgdaFunction{presentation-exact}\AgdaSpace{}%
\AgdaSymbol{=}\AgdaSpace{}%
\AgdaPostulate{proof-elided}\<%
\end{code}

Its proof needs no normal forms --- since both factors are word groups,
$\sem{\_}$ is a splitting rather than a normalisation --- but it does
need, for each of the sixteen rules, that the phase the Weyl cocycle
assigns to the left-hand side equals the correction plus the phase of
the right-hand side.  This reduces to nineteen equations in $\Zp$ about
an explicit phase function on circuits, discharged by Pauli arithmetic,
by a one-wire calculus that reads a circuit as a triple (Pauli, $2
\times 2$ symplectic matrix, phase) with an explicit concatenation law,
and, for \aid{order-SH}, by the number-theoretic fact that the residue
of $(p^2-1)/8$ is $-1/8$.  The extension is provably non-split.  This
is the formal counterpart of the paper's remark before its Theorem~4.10
that a complete rule set for $\langle\omega\rangle$ ``only needs
$\omega^{d} = 1$'', with the one non-trivial correction computed and
checked.
